\section{Conclusion}\label{sec:conclusion}
% When asked about the nature of business cycles, Thomas Sargent\footnote{Hyman Minsky was Thomas Sargent's undergraduate advisor. Whereas Sargent departs methodologically and calls for a quantitative approach to economic research, there is an agreement in their views regarding nature of business cycles.}, a rational expectations pioneer, answered:\footnote{Interview with Thomas Sargent, \emph{The Region}, August 26, 2010. Available at https://www.minneapolisfed.org/article/2010/interview-with-thomas-sargent.}  
% \begin{quote}
% ``[...] economists have been working hard to refine rational expectations theory. [...] An influential example of such work is the 1978 QJE paper by Harrison and Kreps. [...], for policymakers to know whether and how they can moderate bubbles, we need to have well-confirmed quantitative versions of such models up and running."
% \end{quote} 
% This response embraces the idea that ``belief heterogeneity" matters but also calls for quantitative models that link beliefs with the real economy. 

% This paper responds to that call and adapts a standard real business cycle and asset pricing model to fit the narrative that waves of optimism and pessimism drive the business cycle. We argue that a combination of heterogeneity in beliefs, with substantial extrapolation, coupled with an operating leverage channel, can successfully reproduce business cycle, asset pricing and survey data patterns. Extrapolative beliefs and heterogeneity deliver the amplification that is key for this success. Extrapolation adds to risk build up. 

% We foresee that our framework can be extended to study the interaction with other forms of amplification. 
%         Namely, nominal rigidities open the door to amplification through aggregate demand externalities. 
%         Additional financial frictions, such as margin calls or short-selling constraints, can amplify asset prices through fire-sale and pecuniary externalities that impact wealth and risk-taking capacity. 
%         We expect our conclusions to remain relevant in contexts with these additional frictions.


\q{Conclusion}{When asked about the nature of business cycles, Thomas Sargent emphasized that
progress requires quantitative models that discipline deviations from rational
expectations and confront them with the data:\footnote{Interview with Thomas Sargent, \emph{The Region}, August 26, 2010.}
\begin{quote}
"[...] economists have been working hard to refine rational expectations theory. [...] An influential example of such work is the 1978 QJE paper by Harrison and Kreps. [...], for policymakers to know whether and how they can moderate bubbles, we need to have well-confirmed quantitative versions of such models up and running."
\end{quote}

This paper responds to that call. We develop a tractable macro-finance framework
in which heterogeneous and extrapolative beliefs interact with production and labor
demand to generate endogenous movements in risk premia, asset prices, and real
activity.

Quantitatively, the model matches a wide range of empirical patterns in survey expectations,
trading volume, return predictability, and business-cycle dynamics. These results
arise from the interaction between belief dispersion, wealth dynamics, and operating
leverage.

Although our benchmark specification is deliberately parsimonious, with a simple
productivity process and two investor types, the framework is substantially more
general. The Markov structure accommodates richer technology dynamics, stochastic
volatility, alternative belief processes, and additional investor types or frictions.
The simplicity of the baseline clarifies the mechanism; it does not limit the model's
scope.

Our framework provides a disciplined way to embed waves of optimism and pessimism
into a standard macroeconomic environment while remaining flexible enough to
incorporate additional amplification mechanisms. We view this as a step toward
quantitatively credible models in which confidence fluctuations are measurable
forces with testable implications.}